\documentclass[tikz,border=8pt]{standalone}
\usepackage[UTF8]{ctex}
\usepackage{amsmath}
\usetikzlibrary{arrows.meta,decorations.markings}

\definecolor{domainfill}{HTML}{EAF2FF}
\definecolor{outerblue}{HTML}{2563EB}
\definecolor{innerred}{HTML}{DC2626}
\definecolor{cutviolet}{HTML}{7C3AED}
\definecolor{textgray}{HTML}{374151}

\tikzset{
  direction/.style={postaction={decorate},decoration={markings,mark=at position #1 with {\arrow{Stealth[length=3mm]}}}},
  direction/.default=.55
}

\begin{document}
\begin{tikzpicture}[line cap=round,line join=round,font=\small]
  \fill[domainfill] (0,0) ellipse (4.5 and 3.0);
  \fill[white] (0,0) circle (1.05);

  \draw[outerblue,very thick,direction=.25]
    (4.5,0) arc[start angle=0,end angle=360,x radius=4.5,y radius=3.0];

  % The same geometric cut is traversed twice in opposite directions.
  \draw[cutviolet,very thick] (1.05,0) -- (4.50,0);
  \draw[cutviolet,very thick,-{Stealth[length=3mm]}] (4.15,0) -- (3.05,0);
  \draw[cutviolet,very thick,-{Stealth[length=3mm]}] (1.35,0) -- (2.35,0);

  \draw[innerred,very thick,direction=.27]
    (1.05,0) arc[start angle=0,end angle=-360,radius=1.05];

  \fill[innerred] (0,0) circle (2.3pt);
  \node[text=textgray,below=3pt] at (0,0) {奇点};
  \node[text=outerblue] at (-3.25,2.35) {$L$};
  \node[text=innerred,fill=white,inner sep=1.5pt] at (-.92,.82) {$L_{\varepsilon}$};
  \node[text=textgray] at (-2.2,-1.15) {$D_1$};
  \node[text=textgray,align=center] at (2.65,.75) {同一条割线\\往返方向相反};
\end{tikzpicture}
\end{document}
